\scanpage{18}\mspage{-VII- 10 -}
% Faint mark at upper right, possibly showing through from the other side.
\margin{\handwritten{\unsure{XV|XY}}}
\hfill n\textsuperscript{o}383

est l'ensemble des transformés des $u_i$ (resp. $u'_i$) par $G$ (resp. $G'$) (\S~4, n\textsuperscript{o}4, prop.5),
de sorte que $f(S)=S'$. C.Q.F.D.

THEOREME 2.- \uline{Les schémas de Dynkin des systèmes de racines réduits irréductibles
sont à choisir}, \uline{à une homothétie près sur les coefficients}, \uline{parmi les schémas
suivants} :

\diagram{\begin{tikzcd}[ampersand replacement=\&,column sep=1.25em,row sep=.45em]
A_r \& |[inner sep=0pt,label=above:{\scriptsize $1$}]|{\scriptstyle\vert}\arrow[r,dash] \& |[inner sep=0pt,label=above:{\scriptsize $1$}]|{\scriptstyle\vert}\arrow[r,dash] \& |[inner sep=0pt,label=above:{\scriptsize $1$}]|{\scriptstyle\vert}\arrow[r,dash,dashed] \& |[inner sep=0pt,label=above:{\scriptsize $1$}]|{\scriptstyle\vert}\arrow[r,dash] \& |[inner sep=0pt,label=above:{\scriptsize $1$}]|{\scriptstyle\vert}\arrow[r,dash] \& |[inner sep=0pt,label=above:{\scriptsize $1$}]|{\scriptstyle\vert} \& (r\handwritten{\geq}1\text{ noeuds})
\end{tikzcd}}

\diagram{\begin{tikzcd}[ampersand replacement=\&,column sep=1.25em,row sep=.45em]
B_r \& |[inner sep=0pt,label=above:{\scriptsize $1$}]|{\scriptstyle\vert}\arrow[r,dash] \& |[inner sep=0pt,label=above:{\scriptsize $1$}]|{\scriptstyle\vert}\arrow[r,dash] \& |[inner sep=0pt,label=above:{\scriptsize $1$}]|{\scriptstyle\vert}\arrow[r,dash,dashed] \& |[inner sep=0pt,label=above:{\scriptsize $1$}]|{\scriptstyle\vert}\arrow[r,dash] \& |[inner sep=0pt,label=above:{\scriptsize $1$}]|{\scriptstyle\vert}\arrow[r,dash] \& |[inner sep=0pt,label=above:{\scriptsize $(1/2)$}]|{\scriptstyle\vert} \& (r\handwritten{\geq}1\text{ noeuds})
\end{tikzcd}}

\diagram{\begin{tikzcd}[ampersand replacement=\&,column sep=1.25em,row sep=.45em]
C_r \& |[inner sep=0pt,label=above:{\scriptsize $1$}]|{\scriptstyle\vert}\arrow[r,dash] \& |[inner sep=0pt,label=above:{\scriptsize $1$}]|{\scriptstyle\vert}\arrow[r,dash] \& |[inner sep=0pt,label=above:{\scriptsize $1$}]|{\scriptstyle\vert}\arrow[r,dash,dashed] \& |[inner sep=0pt,label=above:{\scriptsize $1$}]|{\scriptstyle\vert}\arrow[r,dash] \& |[inner sep=0pt,label=above:{\scriptsize $1$}]|{\scriptstyle\vert}\arrow[r,dash] \& |[inner sep=0pt,label=above:{\scriptsize $2$}]|{\scriptstyle\vert} \& (r\handwritten{\geq}1\text{ noeuds})
\end{tikzcd}}

\diagram{\begin{tikzcd}[ampersand replacement=\&,column sep=1.25em,row sep=.45em]
{} \& {} \& {} \& {} \& {} \& {} \& |[inner sep=0pt,label=above:{\scriptsize $1$}]|{\scriptstyle\vert}\\
D_r \& |[inner sep=0pt,label=above:{\scriptsize $1$}]|{\scriptstyle\vert}\arrow[r,dash] \& |[inner sep=0pt,label=above:{\scriptsize $1$}]|{\scriptstyle\vert}\arrow[r,dash] \& |[inner sep=0pt,label=above:{\scriptsize $1$}]|{\scriptstyle\vert}\arrow[r,dash,dashed] \& |[inner sep=0pt,label=above:{\scriptsize $1$}]|{\scriptstyle\vert}\arrow[r,dash] \& |[inner sep=0pt,label=above:{\scriptsize $1$}]|{\scriptstyle\vert}\arrow[ur,dash]\arrow[dr,dash] \& {} \& (r\handwritten{\geq}3\text{ noeuds})\\
{} \& {} \& {} \& {} \& {} \& {} \& |[inner sep=0pt,label=below:{\scriptsize $1$}]|{\scriptstyle\vert}
\end{tikzcd}}

\diagram{\begin{tikzcd}[ampersand replacement=\&,column sep=1.25em,row sep=.45em]
E_6,\ E_7,\ E_8 \& |[inner sep=0pt,label=above:{\scriptsize $1$}]|{\scriptstyle\vert}\arrow[r,dash] \& |[inner sep=0pt,label=above:{\scriptsize $1$}]|{\scriptstyle\vert}\arrow[r,dash] \& |[inner sep=0pt,label=above:{\scriptsize $1$}]|{\scriptstyle\vert}\arrow[r,dash,dashed] \& |[inner sep=0pt,label=above:{\scriptsize $1$}]|{\scriptstyle\vert}\arrow[r,dash]\arrow[d,dash] \& |[inner sep=0pt,label=above:{\scriptsize $1$}]|{\scriptstyle\vert}\arrow[r,dash] \& |[inner sep=0pt,label=above:{\scriptsize $1$}]|{\scriptstyle\vert} \& (6,\ 7\text{ ou }8\text{ noeuds})\\
{} \& {} \& {} \& {} \& |[inner sep=0pt,label=right:{\scriptsize $1$}]|{\scriptstyle-}
\end{tikzcd}}

\diagram{\begin{tikzcd}[ampersand replacement=\&,column sep=1.25em,row sep=.45em]
F_4 \& |[inner sep=0pt,label=above:{\scriptsize $1$}]|{\scriptstyle\vert}\arrow[r,dash] \& |[inner sep=0pt,label=above:{\scriptsize $1$}]|{\scriptstyle\vert}\arrow[r,dash] \& |[inner sep=0pt,label=above:{\scriptsize $2$}]|{\scriptstyle\vert}\arrow[r,dash] \& |[inner sep=0pt,label=above:{\scriptsize $2$}]|{\scriptstyle\vert}
\end{tikzcd}}

\diagram{\begin{tikzcd}[ampersand replacement=\&,column sep=1.25em,row sep=.45em]
G_2 \& |[inner sep=0pt,label=above:{\scriptsize $1$}]|{\scriptstyle\vert}\arrow[r,dash] \& |[inner sep=0pt,label=above:{\scriptsize $3$}]|{\scriptstyle\vert}
\end{tikzcd}}

En effet, conservons les notations du début du n\textsuperscript{o}. Soient $X$ (resp. $X'$) le schéma
de Dynkin (resp. Coxeter) de $S$ (resp. $G$). Les possibilités pour $X'$ sont données par
le th.1, les schémas $G_3$ et $G_4$ étant à éliminer puisque l'angle $4\,\handwritten{\pi}/5$ est exclus, de
même que $\underline{D}_2^p$ pour $p\mathrel{\ooalign{\hfil$=$\hfil\cr\hfil$\handwritten{/}$\hfil\cr}}3,4,6$. Quand tous les liens de $X'$ sont d'ordre $3$, tous les
$\lVert u_i\rVert$ sont égaux, d'o\unsure{\=u|\`u}, après homothétie sur $S$, les possibilités $A_r$, $D_r$, $E_6$, $E_7$, $E_8$
pour $X$. D'autre part $\underline{D}_2^3=A_2$, $\underline{D}_2^4=B_2$ ; et, si $X'=\underline{D}_2^6$, $\lVert u_1\rVert^2/\lVert u_2\rVert^2$ est égal à
$3$ ou à $1/3$, d'o\=u, après homothétie, le schéma de Dynkin $X=G_2$. Si $X'=F_4$, les
possibilités pour $X$, après homothétie, sont les suivantes :

\diagram{\begin{tikzcd}[ampersand replacement=\&,column sep=1.25em,row sep=.45em]
|[inner sep=0pt,label=above:{\scriptsize $1$}]|{\scriptstyle\vert}\arrow[r,dash] \& |[inner sep=0pt,label=above:{\scriptsize $1$}]|{\scriptstyle\vert}\arrow[r,dash] \& |[inner sep=0pt,label=above:{\scriptsize $2$}]|{\scriptstyle\vert}\arrow[r,dash] \& |[inner sep=0pt,label=above:{\scriptsize $2$}]|{\scriptstyle\vert}
\end{tikzcd}\qquad \begin{tikzcd}[ampersand replacement=\&,column sep=1.25em,row sep=.45em]
|[inner sep=0pt,label=above:{\scriptsize $2$}]|{\scriptstyle\vert}\arrow[r,dash] \& |[inner sep=0pt,label=above:{\scriptsize $2$}]|{\scriptstyle\vert}\arrow[r,dash] \& |[inner sep=0pt,label=above:{\scriptsize $1$}]|{\scriptstyle\vert}\arrow[r,dash] \& |[inner sep=0pt,label=above:{\scriptsize $1$}]|{\scriptstyle\vert} \& {;}
\end{tikzcd}}

mais ces deux schémas de Dynkin sont isomorphes. Enfin, si $X'=B_r$, on trouve les
seules possibilités $B_r$ et $C_r$ pour $X$.
